\section{Concepts P: Pipeline and Branch Handling}\label{grp:pipe}
{\small\renewcommand{\arraystretch}{1.3}\rowcolors{1}{white}{white}
\begin{longtable}{L{2.3cm}L{4.7cm}L{2.9cm}L{3.0cm}L{2.8cm}}
\toprule
\rowcolor{hdrc}\textbf{Concept} & \textbf{What it fixes $\cdot$ how it works} & \textbf{Best scenario (use when)} & \textbf{Why it beats the alternatives} & \textbf{Weakness $\cdot$ avoid when}\\
\midrule
\endfirsthead
\toprule
\rowcolor{hdrc}\textbf{Concept} & \textbf{What it fixes $\cdot$ how it works} & \textbf{Best scenario (use when)} & \textbf{Why it beats the alternatives} & \textbf{Weakness $\cdot$ avoid when}\\
\midrule
\endhead
\bottomrule
\endfoot
\textbf{\textcolor{cardc}{P1}}\label{card:P1}\newline\textbf{Pipelining}\newline{\footnotesize\textit{Static (design)}} & \textbf{Fixes:} Low instruction throughput. \newline \textbf{How:} Overlap stages (IF, ID, EX, MEM, WB) of successive instructions: throughput up to 1 instr/cycle; latency of one instruction unchanged. & --- & --- & Hazards: structural (same unit), data (RAW/WAR/WAW), control (next PC unknown). Longer pipes increase their impact. \\*
\rowcolor{black!4}\multicolumn{5}{p{\dimexpr\linewidth-2\tabcolsep\relax}}{\textbf{Needs:} Forwarding \cp{P2}, branch handling \cp{P6}--\cp{P12}. \enspace$\cdot$\enspace \textbf{\textcolor{pickc}{Write:}} \textit{Pipelining raises throughput towards 1 instruction/cycle; the [n]-stage pipe loses [x] CPI to hazards, which I remove with [forwarding / prediction / scheduling].}}\\
\midrule
\textbf{\textcolor{cardc}{P2}}\label{card:P2}\newline\textbf{Forwarding / bypassing}\newline{\footnotesize\textit{Static (hardware)}} & \textbf{Fixes:} RAW stalls between dependent ALU operations. \newline \textbf{How:} Route the EX/MEM result straight to the next instruction's input instead of waiting for write-back. & Always in a pipelined core. & Stalling to write-back: removes most RAW bubbles between back-to-back ALU ops with a few muxes. & Cannot remove the load-use bubble (data arrives after MEM); bypass network grows with issue width (Lab 10: extra bypass cycles between unit types). \\*
\rowcolor{black!4}\multicolumn{5}{p{\dimexpr\linewidth-2\tabcolsep\relax}}{\textbf{\textcolor{pickc}{Write:}} \textit{Back-to-back dependent ALU ops [given] would stall until write-back; forwarding passes the EX result directly, leaving only the load-use slot for the compiler.}}\\
\midrule
\textbf{\textcolor{cardc}{P3}}\label{card:P3}\newline\textbf{Compiler fills load/branch delay slots}\newline{\footnotesize\textit{Static}} & \textbf{Fixes:} Load-use and branch bubbles. \newline \textbf{How:} Schedule an independent instruction into the slot after a load or branch. & In-order cores, known latencies, compiler available. & Hardware interlock stalls: same correctness, zero lost cycles when a slot is filled. & Needs independent instructions; deep pipes have too many slots. \\*
\rowcolor{black!4}\multicolumn{5}{p{\dimexpr\linewidth-2\tabcolsep\relax}}{\textbf{Numbers:} FP pipeline latencies (cycles producer$\rightarrow$consumer): FP ALU$\rightarrow$FP ALU 3, FP ALU$\rightarrow$store 2, load$\rightarrow$FP ALU 1, load$\rightarrow$store 0. \enspace$\cdot$\enspace \textbf{\textcolor{pickc}{Write:}} \textit{The compiler moves an independent instruction into the [load/branch] slot, so the bubble does useful work instead of stalling.}}\\
\midrule
\textbf{\textcolor{cardc}{P4}}\label{card:P4}\newline\textbf{Pipelined or replicated FP units}\newline{\footnotesize\textit{Static}} & \textbf{Fixes:} Structural stalls on long-latency units. \newline \textbf{How:} Pipeline the adder/multiplier (initiation interval 1); pipeline or duplicate the divider. & FP-heavy code; divides close together. & Unpipelined divider: initiation interval 25 blocks every next divide for up to 24 cycles. & Out-of-order completion $\Rightarrow$ WAW hazards; register-file write-port conflicts. \newline \textbf{Avoid:} Divides rare: an unpipelined divider saves area and power. \\*
\rowcolor{black!4}\multicolumn{5}{p{\dimexpr\linewidth-2\tabcolsep\relax}}{\textbf{Numbers:} Int ALU lat.\ 0/II 1; load 1/1; FP add 3/1 (4 stages); mult 6/1 (7 stages); divide 24/25. \enspace$\cdot$\enspace \textbf{\textcolor{pickc}{Write:}} \textit{Divides are [frequent and back-to-back]; an unpipelined divider (II 25) stalls each next divide up to 24 cycles, so I pipeline or duplicate it; if divides were rare I would keep it unpipelined to save area.}}\\
\midrule
\textbf{\textcolor{cardc}{P5}}\label{card:P5}\newline\textbf{Deeper pipeline (superpipelining, MIPS R4000)}\newline{\footnotesize\textit{Static}} & \textbf{Fixes:} Clock rate too low. \newline \textbf{How:} Split stages (R4000: 8 stages IF IS RF EX DF DS TC WB; cache access takes 2 stages). & High-clock target, predictable branches, few load-use pairs, strong predictor affordable. & --- & Load delay 2 cycles, branch delay 3 cycles; every stage raises misprediction and load penalty and power. \newline \textbf{Avoid:} Branchy, low-ILP, low-power or small designs: stay shallow. \\*
\rowcolor{black!4}\multicolumn{5}{p{\dimexpr\linewidth-2\tabcolsep\relax}}{\textbf{Needs:} Strong dynamic prediction \cp{P10} + BTB \cp{P11}; small fast L1 \cp{M15}. \enspace$\cdot$\enspace \textbf{\textcolor{pickc}{Write:}} \textit{Deeper pipeline raises $f$ but adds [k] penalty cycles; decide by IC$\times$CPI$\times T_{clk}$, and pair with a [95\%] predictor so $f_{br}\times$miss$\times$penalty stays small.}}\\
\midrule
\textbf{\textcolor{cardc}{P6}}\label{card:P6}\newline\textbf{Early branch resolution (compare + target in ID)}\newline{\footnotesize\textit{Static}} & \textbf{Fixes:} Branch penalty in a short pipeline. \newline \textbf{How:} Move compare and target calculation into ID. & Cheap 5-stage in-order cores. & Resolving in EX (2 stalls) or MEM (3 stalls). & Comparator and adder in ID; extra stall for instructions producing the branch operand. \\*
\rowcolor{black!4}\multicolumn{5}{p{\dimexpr\linewidth-2\tabcolsep\relax}}{\textbf{Numbers:} Stall per branch: resolve in MEM 3, EX 2, \textbf{ID 1}. Side effect: an ALU producer of the branch operand needs +1 stall, a load producer +2. \enspace$\cdot$\enspace \textbf{\textcolor{pickc}{Write:}} \textit{Moving the compare and target add into ID cuts the penalty from [3] to 1 cycle per branch, i.e.\ $[f_{br}]\times2$ CPI saved, with almost no hardware.}}\\
\midrule
\textbf{\textcolor{cardc}{P7}}\label{card:P7}\newline\textbf{Static prediction: not-taken, BTFNT, profile}\newline{\footnotesize\textit{Static}} & \textbf{Fixes:} Branch penalty without predictor hardware. \newline \textbf{How:} \emph{Not taken}: keep fetching fall-through, flush if taken. \emph{BTFNT}: backward branches (loops) taken, forward not taken. \emph{Profile}: compiler hint from earlier runs. & Cost-sensitive, embedded, fixed firmware, loop-dominated code, real-time. & Dynamic predictor when behaviour is known offline: zero area, deterministic. & \textbf{Avoid:} Data-dependent branches, deep pipe; predict-not-taken alone fails on loops. ``Predict taken'' is useless unless the target is known early (BTB). \\*
\rowcolor{black!4}\multicolumn{5}{p{\dimexpr\linewidth-2\tabcolsep\relax}}{\textbf{\textcolor{pickc}{Write:}} \textit{The firmware is fixed and loop-dominated, so BTFNT predicts the back-edges correctly with zero predictor area; a dynamic predictor's area is not justified in a price-driven design.}}\\
\midrule
\textbf{\textcolor{cardc}{P8}}\label{card:P8}\newline\textbf{Delayed branch; cancelling (branch-likely)}\newline{\footnotesize\textit{Static}} & \textbf{Fixes:} The 1 slot after a branch. \newline \textbf{How:} Slot instruction always executes. Fill it from \emph{before} (always safe), from the \emph{target} (branch usually taken, e.g.\ loops; must be harmless otherwise), or from \emph{fall-through} (usually not taken). Cancelling branch nullifies the slot if the prediction is wrong (MIPS branch-likely = cancel if not taken: ideal for loop back-edges). & --- & Plain stall: about half the slots do useful work; cancelling lets the compiler fill from the target without being conservative. & \textbf{Avoid:} Deep pipes (several slots, hard to fill). \\*
\rowcolor{black!4}\multicolumn{5}{p{\dimexpr\linewidth-2\tabcolsep\relax}}{\textbf{Numbers:} Compiler fills $\approx$50\% of single slots usefully. \enspace$\cdot$\enspace \textbf{\textcolor{pickc}{Write:}} \textit{With one delay slot and loop back-edges [given], filling from the target with a branch-likely instruction turns the bubble into useful work on every taken iteration.}}\\
\midrule
\textbf{\textcolor{cardc}{P9}}\label{card:P9}\newline\textbf{2-bit BHT (saturating counters)}\newline{\footnotesize\textit{Dynamic}} & \textbf{Fixes:} Mispredictions of loops and biased branches. \newline \textbf{How:} Table indexed by low PC bits, no tags (a hint); prediction flips only after two wrong outcomes. & Small/cheap cores, loop code. & 1-bit: a 10-iteration loop mispredicts twice (exit and re-entry) $\Rightarrow$ 80\%; 2-bit only once. $n$-bit gives little more. & \textbf{Avoid:} Branches correlated with other branches (use \cp{P10}); deep/wide pipelines where 90\% is not enough. \\*
\rowcolor{black!4}\multicolumn{5}{p{\dimexpr\linewidth-2\tabcolsep\relax}}{\textbf{Numbers:} 4K-entry 2-bit BHT: 82--99\% accuracy (SPEC89), nearly an infinite table. Beyond 4K entries make it \emph{smarter}, not bigger. \enspace$\cdot$\enspace \textbf{\textcolor{pickc}{Write:}} \textit{A 4K-entry 2-bit BHT (1\,KB, no tags) gives 82--99\% accuracy; loops mispredict once per exit instead of twice.}}\\
\midrule
\textbf{\textcolor{cardc}{P10}}\label{card:P10}\newline\textbf{Correlating / two-level / gshare predictor}\newline{\footnotesize\textit{Dynamic}} & \textbf{Fixes:} Branches whose outcome depends on other recent branches (if-chains, control code). \newline \textbf{How:} Global history of last $m$ outcomes selects among $2^m$ counter tables: $(m,n)$ predictor; two-level = history register + pattern table; \textbf{gshare} indexes the table with PC XOR history (better table use\bs). & Branchy integer code, interpreters, compilers, deep or wide pipelines. & A small $(2,2)$ beats a same-size 2-bit table and often an \emph{unlimited} 2-bit table. & More area/energy; warm-up; aliasing. \\*
\rowcolor{black!4}\multicolumn{5}{p{\dimexpr\linewidth-2\tabcolsep\relax}}{\textbf{Numbers:} $(2,2)$ with 64 entries: 6-bit index $=$ 2 history + 4 PC bits. \enspace$\cdot$\enspace \textbf{\textcolor{pickc}{Write:}} \textit{With [20\%] branches and a [3]-cycle penalty, raising accuracy from 90\% to 95\% cuts branch CPI from 0.06 to 0.03; correlated if-chains in [workload] justify a two-level predictor.}}\\
\midrule
\textbf{\textcolor{cardc}{P11}}\label{card:P11}\newline\textbf{Branch target buffer (BTB), target buffer, fetch unit}\newline{\footnotesize\textit{Dynamic}} & \textbf{Fixes:} Target address unknown in IF (bubble even if direction is right). \newline \textbf{How:} Cache of taken branches: tag + predicted target (+ bits), read in IF with the PC. Tag must match. Target-instruction buffer stores the instruction itself $\Rightarrow$ branch folding (0-cycle jumps). Integrated fetch unit = prediction + prefetch + buffering. & Any predicted-taken scheme; wide issue; deep pipe. & Direction-only predictor: removes the target bubble, so a correctly predicted taken branch costs 0 cycles. & Area; misprediction: kill, refetch, update. \\*
\rowcolor{black!4}\multicolumn{5}{p{\dimexpr\linewidth-2\tabcolsep\relax}}{\textbf{Needs:} Accurate direction predictor (\cp{P9}/\cp{P10}); recovery on mispredict. \enspace$\cdot$\enspace \textbf{\textcolor{pickc}{Write:}} \textit{Taken branches every [k] instructions [given] each lose a fetch cycle without the target; a BTB supplies it in IF, so predicted-taken branches cost nothing.}}\\
\midrule
\textbf{\textcolor{cardc}{P12}}\label{card:P12}\newline\textbf{Return-address stack}\newline{\footnotesize\textit{Dynamic}} & \textbf{Fixes:} Returns/indirect jumps whose target changes at run time. \newline \textbf{How:} Push on call, pop to predict the return. & Call-heavy code (OO, recursion, libraries). & BTB for returns: a return's target changes with the call site; the stack is exact for nested calls. & Overflow on very deep recursion\bs. \\*
\midrule
\textbf{\textcolor{cardc}{P13}}\label{card:P13}\newline\textbf{Predication / branch-free code\bs}\newline{\footnotesize\textit{Static}} & \textbf{Fixes:} Unpredictable short branches; SIMT divergence. \newline \textbf{How:} Turn \texttt{if} into conditional moves/masked ops: both sides computed, result selected. & Short, data-dependent \texttt{if} bodies; GPUs/SIMD (masks). & Prediction when outcome is random (50/50): no flush penalty. & Both sides always execute (energy, issue slots); needs predicated/select instructions. \newline \textbf{Avoid:} Long bodies: you always pay both sides. \\*
\rowcolor{black!4}\multicolumn{5}{p{\dimexpr\linewidth-2\tabcolsep\relax}}{\textbf{\textcolor{pickc}{Write:}} \textit{The branch is data-dependent and close to 50/50 [given], so no predictor helps; converting the short \texttt{if} to a conditional select removes the [n]-cycle flush entirely.}}\\
\end{longtable}}

\begin{ex}[Branch numbers to quote: penalty 3, 20\% branches, 60\% taken, base CPI 1]
Stall always $+0.60$ (CPI 1.60); predict-not-taken $0.2\times0.6\times3=+0.36$ (1.36); 2-bit+BTB 90\% $+0.06$ (1.06); two-level 95\% $+0.03$ (1.03). Speedup of dynamic over stall $\approx1.5\times$; two-level over 2-bit only $\approx1.03\times$ (diminishing returns: justify only for deep/wide pipes or high branch frequency). Even a 1-cycle penalty costs 10--30\% (slide).
\end{ex}

\small
\begin{tabularx}{\textwidth}{L{5.0cm}L{4.4cm}X}
\toprule
\tblhead{\textbf{Scenario} & \textbf{Choose} & \textbf{Because}}\\
\midrule
Small/cheap core, loop code & 2-bit BHT \cp{P9} & Fixes the loop double-miss with 2 bits/entry, no tags \\
Fixed firmware, no predictor area & BTFNT/profile + delay slot \cp{P7} \cp{P8} & Behaviour known offline; zero hardware \\
Control-heavy integer, correlated if-chains & $(2,2)$/two-level/gshare \cp{P10} & Uses other branches' outcomes; small table beats a bigger 2-bit one \\
Deep or wide pipeline & Two-level + BTB + RAS + speculation \cp{P10} \cp{P11} \cp{P12} \cp{I5} & Penalty large; wide issue needs a branch per cycle \\
Many calls/returns, indirect jumps & Return stack + target-predicting BTB \cp{P12} \cp{P11} & Run-time targets \\
Random 50/50 short branch & Predication \cp{P13} & No predictor can learn randomness \\
\bottomrule
\end{tabularx}
\normalsize
